\bookchapter{Numerical derivatives}\label{ch:numerical}

\section*{Explanation}
A computer can approximate a derivative without any rules, using a very small step $h$ (for example $h=10^{-5}$):
\begin{align*}
f'(x)&\approx\frac{f(x+h)-f(x)}{h} &&\text{(forward difference)},\\
f'(x)&\approx\frac{f(x+h)-f(x-h)}{2h} &&\text{(central difference)}.
\end{align*}
The central difference is usually more accurate. This is useful for checking derivatives computed by hand (\emph{gradient checking}).

\begin{example}
$f(x)=x^2$ at $x=3$ with $h=0.1$:
\[
\text{forward: }\frac{3.1^2-3^2}{0.1}=\frac{9.61-9}{0.1}=6.1,\qquad
\text{central: }\frac{3.1^2-2.9^2}{0.2}=\frac{9.61-8.41}{0.2}=6.
\]
The exact value is $f'(3)=6$.
\end{example}

\begin{example}[Python]
\begin{verbatim}
def derivative(f, x, h=1e-5):
    return (f(x + h) - f(x - h)) / (2 * h)

print(derivative(lambda x: x**2, 3))   # approx. 6.0
\end{verbatim}
\end{example}

\section*{Practice problems}
\begin{enumerate}[label=\textbf{\thechapter.\arabic*}]
  \item Approximate $f'(2)$ for $f(x)=x^3$ with the central difference and $h=0.1$. Compare with the exact value.
  \item Approximate $f'(1)$ for $f(x)=x^2$ with the forward difference and $h=0.5$. Compare with the exact value.
  \item \texttt{[code]} Use the function \texttt{derivative} above to approximate $\sin'(0)$, $\sin'(\pi/2)$ and $\exp'(1)$. Compare with $\cos$ and $e$.
  \item \texttt{[code]} Implement the sigmoid, compute its numerical derivative at $x=0$ and compare to $\tfrac14$.
  \item[\textbf{9.5}$\star$] \texttt{[code]} Experiment: for $f(x)=e^x$ at $x=1$ compute the error of the forward difference for $h=10^{-1},10^{-3},10^{-5},10^{-8},10^{-12}$. What do you observe for very small $h$? (Computers store numbers with limited precision.)
\end{enumerate}

% =====================================================
